\documentclass[10pt,a4paper]{amsart}
\usepackage[margin=19mm]{geometry}
\usepackage{amsmath,amssymb,mathtools}
\usepackage[scr=rsfs,cal=euler]{mathalfa}
\usepackage{xfrac}
\usepackage[unicode,hidelinks,bookmarks=false]{hyperref}
\newcommand{\ie}{i.e.\ }
\newcommand{\cf}{cf.\ }
\newcommand{\resp}{respectively\ }
\newcommand{\Subsection}{\subsection}
\providecommand{\nobreakdash}{\nobreak\hbox{-}}
\setlength{\emergencystretch}{2em}
\multlinegap0pt
\usepackage{bm}
\usepackage{bbm}
\usepackage{dsfont}
\usepackage{xspace}
\usepackage{cancel}
\usepackage{mathtools}
\usepackage{amsthm}
\makeatletter
\@ifundefined{thm}{\newtheorem{thm}{Thm}}{}
\@ifundefined{prop}{\newtheorem{prop}[thm]{Proposition}}{}
\makeatother
\title[Pointwise character bounds for $\mathrm{SU}(3)$]{Pointwise character bounds for $\mathrm{SU}(3)$}
\author{Yunfeng Zhang}
\date{Source: arXiv:2604.17589v1 (CC BY 4.0)}
\begin{document}
\maketitle

\noindent\emph{Source and licence.} Pointwise character bounds for $\mathrm{SU}(3)$. Version arXiv:2604.17589v1 (CC BY 4.0), distributed under the Creative Commons Attribution 4.0 International licence, as recorded in the arXiv licence metadata for this exact version and shown on the abstract page https://arxiv.org/abs/2604.17589v1. Full rights notice: licenses/arxiv-2604.17589-CC-BY-4.0.txt.\par\smallskip

\noindent\textit{Editorial scope.} Counted unit: the Prop whose source label is \texttt{prop: I} in the article (source0.tex lines 1037--1052, source label \texttt{prop: I}), delivered together with the article's own complete original proof of it (source0.tex lines 1053--1198, one \texttt{proof} environment of 146 lines). No mathematics is written, repaired or re-derived: statement and proof are the article's own text line for line. Required context retained uncounted: eq: ABC integral 2 (lines 1030--1035). Every definition, display and label of those lines is retained unchanged. Omitted from this package and not redistributed: 1--1029, 1036--1036, 1199--1229. Nothing inside a retained excerpt is omitted or altered: every retained body is delivered exactly as the article's lines are written, so no omission receipt of any kind accompanies this package. Citations: the retained excerpts carry no citation command at all, so no bibliography entry of the article is transcribed and no reference is redistributed. Reference adaptation: none was needed and none was made; every reference inside the delivered unit resolves inside it, so no unresolved reference or citation is produced. Printed slips kept literal: no printed word, symbol, hypothesis or number of the counted statement or of its proof was altered. Layout and class: the article's own class is \texttt{amsart} with packages including amsmath, amssymb, hyperref, enumitem, mathrsfs, rank-2-roots, tikz, etoolbox; the wrapper is the lane's pinned \texttt{amsart} editorial wrapper at 10pt on A4, which restates the theorem-like environments the retained text needs and adds the article's own macro definitions that the retained text needs (none), with extra wrapper packages (bm, bbm, dsfont, xspace, cancel, mathtools). The delivered numbering is the unit's own. Assets: no retained span contains a figure environment, a graphics-inclusion command, a PDF-page inclusion, a table or a TikZ picture; no image file is copied into this package, the package declares no asset dependency and no image file is redistributed with it. Licence: version arXiv:2604.17589v1 of this article is distributed under the Creative Commons Attribution 4.0 International licence, as recorded in the arXiv licence metadata for this exact version and shown on the abstract page https://arxiv.org/abs/2604.17589v1. A full rights notice is delivered at licenses/arxiv-2604.17589-CC-BY-4.0.txt. Characters: every retained excerpt is pure ASCII, so the delivered engine, XeLaTeX with its default Latin Modern text font, reports no missing character.
\begin{align}
\mathcal{I}(s):=&\int_{A_0} \prod_{\alpha\in\widetilde{\Phi}}
   \frac{\|\langle\alpha,H\rangle\|^2}{(1+|\langle s(\mu+\rho),\alpha\rangle|\cdot\|\langle\alpha, H\rangle\|)^{p}}
   \,dH \label{eq: ABC integral} \\
   \lesssim& \int_{A_0}\frac{t_1^2t_2^2(t_1+t_2)^2}{(1+\widetilde{a}t_1)^p(1+\widetilde{b}t_2)^p(1+\widetilde{c}(t_1+t_2))^p}\,dt_1\,dt_2. \label{eq: ABC integral 2}
\end{align}

\begin{prop}\label{prop: I}
Let $(a,b,c)$ be the relabeling of the set 
$\{ |\langle \mu+\rho, \alpha \rangle| : \alpha \in \widetilde{\Phi} \}$ such that $a\ge b\ge c$. Then for all $s\in W$, we have 
    \[
\mathcal I(s)\lesssim
\begin{cases}
a^{-p}b^{-p}c^{-p}, & 0<p<\frac83,\\[1mm]
a^{-p}b^{-p}c^{-p}\log(2+c), & p=\frac83,\\[1mm]
a^{-p}b^{-p}c^{2p-8}, & \frac83<p<3,\\[1mm]
a^{-3}b^{-3}c^{-2}\log\left(2+\frac ac\right), & p=3, \\[1mm]
a^{-3}b^{-p}c^{p-5}, & 3<p<5,\\[1mm]
a^{-3}b^{-5}\log\left(2+\frac bc\right), & p=5,\\[1mm]
a^{-3}b^{-5}, & p>5.
\end{cases}
\]
\end{prop}

\begin{proof}
As $s$ varies over $W$, the triple $(\widetilde{a},\widetilde{b},\widetilde{c})$ runs over all permutations of the fixed set
$\{ |\langle \mu+\rho, \alpha \rangle| : \alpha \in \widetilde{\Phi} \}$. Observe that the integral in \eqref{eq: ABC integral 2} is maximized when
\[
(\widetilde a,\widetilde b,\widetilde c)=(a,b,c)
\quad\text{or}\quad
(b,a,c).
\]
Thus it suffices to get the desired estimate for the integral in \eqref{eq: ABC integral 2} when $(\widetilde a,\widetilde b,\widetilde c)=(a,b,c)$. 

Splitting \(A_0\) into \(t_1\ge t_2\) and \(t_1\le t_2\), we may further assume 
\[t_1\le t_2,\]
since this region yields the larger contribution under the assumption $a\geq b$. Then $t_1+t_2\approx t_2$. 

Let $M=\frac{4\pi}{3}$. We have 
\begin{align*}
    \mathcal{I}=\mathcal{I}(s)&\lesssim \int_{0\leq t_1\leq t_2\leq M} 
    \frac{t_1^2t_2^4}{(1+at_1)^p(1+bt_2)^p(1+ct_2)^p}
    \,dt_1\,dt_2\\
    &=\int_0^M \frac{t_2^4}{(1+bt_2)^p(1+ct_2)^p}
\left(\int_0^{t_2}\frac{t_1^2}{(1+at_1)^p}\,dt_1\right)\,dt_2.
\end{align*}

\underline{The case of \(0<p<3\).} We have
\[
\int_0^{t_2}\frac{t_1^2}{(1+at_1)^p}\,dt_1
\lesssim
\begin{cases}
t_2^3, & 0\le t_2\le a^{-1},\\[1mm]
a^{-p}t_2^{3-p}, & t_2\ge a^{-1}.
\end{cases}
\]
Recall that $a \ge b \ge c > 0$. We split the $t_2$-integral at 
$t_2 = a^{-1}, b^{-1}, c^{-1}$ to obtain 
\begin{align*}
\mathcal I
\lesssim&
\int_0^{a^{-1}} t_2^7\,dt_2
+a^{-p}\int_{a^{-1}}^{b^{-1}} t_2^{7-p}\,dt_2\\
& +a^{-p}b^{-p}\int_{b^{-1}}^{c^{-1}} t_2^{7-2p}\,dt_2
+a^{-p}b^{-p}c^{-p}\int_{c^{-1}}^M t_2^{7-3p}\,dt_2.
\end{align*}
For the first three terms, we have 
\begin{align*}
    & \int_0^{a^{-1}} t_2^7\,dt_2
+a^{-p}\int_{a^{-1}}^{b^{-1}} t_2^{7-p}\,dt_2
+a^{-p}b^{-p}\int_{b^{-1}}^{c^{-1}} t_2^{7-2p}\,dt_2\\
\lesssim & a^{-8}+a^{-p}b^{p-8}+a^{-p}b^{-p}c^{2p-8}.
\end{align*}
For the last term, we have 
$$a^{-p}b^{-p}c^{-p}\int_{c^{-1}}^M t_2^{7-3p}\,dt_2\lesssim \begin{cases}
a^{-p}b^{-p}c^{-p}, & 0<p<\frac83,\\[1mm]
a^{-p}b^{-p}c^{-p}\log(2+c), & p=\frac83,\\[1mm]
a^{-p}b^{-p}c^{2p-8}, & \frac83<p<3.
\end{cases}$$
Using $a\geq b\geq c$,  we obtain
\[
\mathcal I\lesssim
\begin{cases}
a^{-p}b^{-p}c^{-p}, & 0<p<\frac83,\\[1mm]
a^{-p}b^{-p}c^{-p}\log(2+c), & p=\frac83,\\[1mm]
a^{-p}b^{-p}c^{2p-8}, & \frac83<p<3.
\end{cases}
\]

\underline{The case of $p=3$.} We have
\begin{align*}
\int_0^{t_2}\frac{t_1^2}{(1+at_1)^3}\,dt_1
\lesssim
\begin{cases}
t_2^3, & 0\le t_2\le a^{-1},\\[4pt]
a^{-3}\log(2+at_2), & t_2\ge a^{-1}.
\end{cases}
\end{align*}
Then we have 
\begin{align*}
    \mathcal{I}
    \lesssim &\int_0^{a^{-1}}t_2^7\, dt_2
    +a^{-3}\int_{a^{-1}}^{b^{-1}}t_2^4\log(2+at_2)\, dt_2\\
    &+a^{-3}b^{-3}\int_{b^{-1}}^{c^{-1}}t_2\log(2+at_2)\, dt_2
    +a^{-3}b^{-3}c^{-3}\int_{c^{-1}}^{M}t_2^{-2}\log(2+at_2)\, dt_2.
\end{align*}
The last integral above may be estimated by 
\begin{align*}
 \int_{c^{-1}}^{M}t_2^{-2}\log(2+at_2)\,dt_2&
 \lesssim \int_{c^{-1}}^{M}t_2^{-2}\left[\log\left(2+\frac ac\right)+\log(ct_2)\right] \,dt_2\\
 &\lesssim c\log\left(2+\frac ac\right),
\end{align*}
so that 
\begin{align*}
    \mathcal{I}
    \lesssim &  a^{-8}+a^{-3}b^{-5}\log\left(2+\frac ab\right)
    +a^{-3}b^{-3}c^{-2}\log\left(2+\frac ac\right).
\end{align*}
Since $a\ge b\ge c$, the last term dominates, and hence
\begin{align*}
\mathcal I
\lesssim a^{-3}b^{-3}c^{-2}\log\!\left(2+\frac ac\right).
\end{align*}

\underline{The case of \(p>3\).} We have 
\[
\int_0^{t_2}\frac{t_1^2}{(1+at_1)^p}\,dt_1
\lesssim
\begin{cases}
t_2^3, & 0\le t_2\le a^{-1},\\[1mm]
a^{-3}, & t_2\ge a^{-1}.
\end{cases}
\]
Hence
\begin{align*}
\mathcal I
\lesssim & 
\int_0^{a^{-1}} t_2^7\,dt_2
+a^{-3}\int_{a^{-1}}^{b^{-1}} t_2^4\,dt_2
 \\
&+a^{-3}b^{-p}\int_{b^{-1}}^{c^{-1}} t_2^{4-p}\,dt_2+a^{-3}b^{-p}c^{-p}\int_{c^{-1}}^M t_2^{4-2p}\,dt_2.
\end{align*}
For the third term, we have 
\[
a^{-3}b^{-p}\int_{b^{-1}}^{c^{-1}} t_2^{4-p}\,dt_2
\lesssim
\begin{cases}
a^{-3}b^{-p}c^{p-5}, & 3<p<5,\\[1mm]
a^{-3}b^{-5}\log\left(2+\frac bc\right), & p=5,\\[1mm]
a^{-3}b^{-5}, & p>5.
\end{cases}
\]
For the remaining terms, we have 
\begin{align*}
    & \int_0^{a^{-1}} t_2^7\,dt_2
+a^{-3}\int_{a^{-1}}^{b^{-1}} t_2^4\,dt_2
 +a^{-3}b^{-p}c^{-p}\int_{c^{-1}}^M t_2^{4-2p}\,dt_2 \\ 
 \lesssim & a^{-8}+a^{-3}b^{-5}+a^{-3}b^{-p}c^{p-5}. 
\end{align*}
Using $a\geq b\geq c$, we obtain
\[
\mathcal I\lesssim
\begin{cases}
a^{-3}b^{-p}c^{p-5}, & 3<p<5,\\[1mm]
a^{-3}b^{-5}\log\left(2+\frac bc\right), & p=5,\\[1mm]
a^{-3}b^{-5}, & p>5.
\end{cases}
\]
The proof is completed. 
\end{proof}

\end{document}
